% =====================================================================
\chapter{Диагностика неполадок}
% =====================================================================
\chapterintro
  {типичные проблемы с пультом, привязкой, телеметрией и полётным контроллером и
   порядок их поиска}
  {терпение и эта таблица}

\section{Общий порядок поиска}

\begin{enumerate}
  \item \textbf{Где пропадает сигнал?} Идите по конвейеру (рис.~\ref{fig:pipeline}): стик
        двигается на экране калибровки? вход меняется? канал в \menu{Outputs} меняется?
        приёмник получает (светодиод, телеметрия)? контроллер видит (вкладка Receiver)?
  \item \textbf{Что менялось последним?} Прошивка, модель, приёмник, проводка.
  \item \textbf{Упростите}: сделайте новую «чистую» модель и проверьте на ней.
\end{enumerate}

\section{Таблица проблем}

\begin{small}
\renewcommand{\arraystretch}{1.2}
\begin{longtable}{@{}p{48mm}p{112mm}@{}}
\toprule
\textbf{Симптом} & \textbf{Причина и решение} \\
\midrule
\endhead
Пульт не включается & разряжены или неправильно вставлены элементы; держите кнопку питания
  дольше; для 18650 проверьте, достают ли элементы до контактов. \\
При включении «SD card warning» / «Storage warning» & версия содержимого SD-карты не
  совпадает с прошивкой: обновите SD через Buddy (глава~\ref{ch:firmware}). \\
Застряли на «Throttle warning» / «Switch warning» & переведите газ и переключатели в
  стартовые положения; проверьте, не требует ли модель положения, которое физически
  невозможно (переставленный тумблер) --- исправьте в \menu{Preflight checks}. \\
Стик в центре показывает не 0, по краям не 100 & повторите калибровку (глава~\ref{ch:radio}). \\
Средняя позиция трёхпозиционного тумблера не работает & в \menu{Hardware} тип указан как
  \menu{2POS}; установите \menu{3POS}. \\
В скрипте ExpressLRS «Loading…» или \code{--} & \menu{Internal RF} выключен или стоит не
  \menu{CRSF}; неверный тип модуля в \menu{Hardware}; слишком высокая \menu{Baudrate} для
  модуля --- попробуйте 400\,000. \\
ELRS не привязывается & разные фразы привязки; разные мажорные версии ELRS на передатчике и
  приёмнике; разный частотный диапазон (2,4\,ГГц vs 900\,МГц); включён Model Match с другим
  номером приёмника. \\
Приёмник FrSky не привязывается & подтип (D8/D16/LBT/FCC) или версия ACCST (v1/v2) не
  совпадает с прошивкой приёмника; попробуйте \menu{RF Freq. fine tune}. \\
Связь есть, но каналы «перепутаны» & порядок каналов пульта (AETR) не совпадает с
  \menu{Channel map} контроллера. \\
Коптер не армится & \sw{CH5} не доходит до диапазона ARM в Betaflight; газ не на минимуме;
  коптер не на ровной поверхности; причина видна в Betaflight (\code{arming disable flags}
  в CLI командой \code{status}). \\
Телеметрия не приходит & включите телеметрию (\menu{Telem Ratio} не \menu{Off}); для FrSky ---
  подтип с телеметрией; повторите \menu{Discover new sensors}. \\
Частые «Telemetry lost» вблизи & антенна приёмника повреждена или экранирована карбоном;
  слишком высокая мощность вплотную перегружает приёмник --- отойдите на 2--3\,м. \\
Нет голоса & \menu{Sound mode} запрещает голос; громкость 0; нет папки \code{SOUNDS/<язык>}
  или язык в \menu{Voice language} другой. \\
Не видно пульт на компьютере & кабель только для зарядки; режим USB не выбран; для DFU на
  Windows --- драйвер WinUSB (Zadig). \\
Lua-скрипт не запускается & не та папка; слишком длинное имя файла; синтаксическая ошибка;
  скрипт написан для цветного экрана. \\
Boxer: вентилятор шумит постоянно & высокая мощность ELRS; поднимите \menu{Fan Thresh}
  или снизьте мощность. \\
Быстро садится батарея пульта & мощность модуля 500--1000\,мВт, яркая подсветка,
  фоновая музыка; уменьшите мощность и включите динамическую. \\
\bottomrule
\end{longtable}
\end{small}

\section{Полезные привычки}

\begin{itemize}
  \item Раз в месяц --- резервная копия SD-карты.
  \item Перед сезоном --- проверка failsafe и калибровки.
  \item Одна модель в пульте --- одна модель в жизни. Имена моделей подписаны и на моделях.
  \item Заряд пульта --- перед каждым выездом, а не «вроде ещё хватит».
\end{itemize}
